% GENERATED FILE. Edit canonical lessons, then run npm run generate:books.
% canonical-source-hash: fnv1a64:05b9a4d24aa7a0f0
% canonical-lessons: GU-C481-kapavu, GU-C481-rujhavu, GU-C481-ghavavu, GU-C481-jharvu, GU-C481-cholavu

\chapter{Cuts and Scrapes}
\label{ch:gu-cuts-and-scrapes}
\hlchaptermodality{481}

\begin{chapteropening}
\textbf{By the end of this chapter:} \emph{I can say how I got hurt and how it healed.}

Complete the last lesson of chapter 481: I can say how I got hurt and how it healed.
\end{chapteropening}

\section[\texorpdfstring{\gu{કપાવું} (kapāvũ)}{કપાવું (kapāvũ)}]{\gu{કપાવું} (kapāvũ) — to get cut}

\label{lesson:GU-C481-kapavu}



\begin{quote}
Before the new one: say the Gujarati for to rot, then the Gujarati for to puff up.
\end{quote}

\subsection*{You'll want to know: \gu{કપાવું}}
\textbf{\gu{કપાવું}} — \emph{kapāvũ} — \textquotedblleft{}to get cut\textquotedblright{}.

\begin{quote}
\gu{છરીથી} \gu{મારી} \gu{આંગળી} \gu{કપાઈ} \gu{ગઈ}.
\end{quote}

\begin{quote}
\emph{charīthī mārī ā̃gaḷī kapāī gaī.}
\end{quote}

\begin{quote}
I cut my finger on the knife.
\end{quote}

Past (it …): \gu{તે} \gu{આંગળી} \gu{કપાઈ} \gu{ગઈ}, \emph{te ā̃gaḷī kapāī gaī}.

\begin{cousinweb}
\textbf{\gu{કપાવું}} (\emph{kapāvũ}) is the intransitive partner of \textbf{\gu{કાપવું}} (\emph{kāpvũ}), \textquotedblleft{}to cut\textquotedblright{}, which goes back to Sanskrit \emph{kalpayati}, \textquotedblleft{}trims, cuts\textquotedblright{} (CDIAL 2944).
\end{cousinweb}

\begin{grammarlens}[title={What you've built}]
One more word for describing what you lived and felt.
\end{grammarlens}

\subsection*{Your turn}
\begin{itemize}
\raggedright
  \item \emph{Say it:} \emph{kapāvũ}
  \item \emph{Say it:} \emph{kapāvũ}, once more
  \item \emph{Say it:} \emph{charīthī mārī ā̃gaḷī kapāī gaī.} — I cut my finger on the knife.
  \item \emph{Recall:} say the Gujarati for to rot, then the Gujarati for to puff up, then say \emph{kapāvũ} again
\end{itemize}

\begin{tcolorbox}[breakable,colback=teal!4,colframe=teal!35!black,title={Before you move on}]
What does \gu{કપાવું} mean? (\textquotedblleft{}To get cut\textquotedblright{}.) Say it once more.
\end{tcolorbox}
\section[\texorpdfstring{\gu{રુઝાવું} (rujhāvũ)}{રુઝાવું (rujhāvũ)}]{\gu{રુઝાવું} (rujhāvũ) — to heal (of a wound)}

\label{lesson:GU-C481-rujhavu}



\begin{quote}
Before the new one: say the Gujarati for to puff up, then the Gujarati for to get cut.
\end{quote}

\subsection*{You'll want to know: \gu{રુઝાવું}}
\textbf{\gu{રુઝાવું}} — \emph{rujhāvũ} — \textquotedblleft{}to heal (of a wound)\textquotedblright{}.

\begin{quote}
\gu{આ} \gu{ઘા} \gu{એક} \gu{અઠવાડિયામાં} \gu{રુઝાઈ} \gu{જશે}.
\end{quote}

\begin{quote}
\emph{ā ghā ek aṭhvāḍiyāmā̃ rujhāī jaśe.}
\end{quote}

\begin{quote}
This wound will heal in a week.
\end{quote}

Past (it …): \gu{તે} \gu{ઘા} \gu{રુઝાઈ} \gu{ગયો}, \emph{te ghā rujhāī gayo}.

\begin{cousinweb}
\textbf{\gu{રુઝાવું}} (\emph{rujhāvũ}) is inherited from Sanskrit \emph{rudhyatē}, \textquotedblleft{}is stopped, is obstructed\textquotedblright{}: a wound that closes. CDIAL lists the Gujarati verb, \textquotedblleft{}to be healed, of a wound\textquotedblright{} (CDIAL 10779).
\end{cousinweb}

\begin{grammarlens}[title={What you've built}]
One more word for describing what you lived and felt.
\end{grammarlens}

\subsection*{Your turn}
\begin{itemize}
\raggedright
  \item \emph{Say it:} \emph{rujhāvũ}
  \item \emph{Say it:} \emph{rujhāvũ}, once more
  \item \emph{Say it:} \emph{ā ghā ek aṭhvāḍiyāmā̃ rujhāī jaśe.} — This wound will heal in a week.
  \item \emph{Recall:} say the Gujarati for to puff up, then the Gujarati for to get cut, then say \emph{rujhāvũ} again
\end{itemize}

\begin{tcolorbox}[breakable,colback=teal!4,colframe=teal!35!black,title={Before you move on}]
What does \gu{રુઝાવું} mean? (\textquotedblleft{}To heal (of a wound)\textquotedblright{}.) Say it once more.
\end{tcolorbox}
\section[\texorpdfstring{\gu{ઘવાવું} (ghavāvũ)}{ઘવાવું (ghavāvũ)}]{\gu{ઘવાવું} (ghavāvũ) — to be wounded}

\label{lesson:GU-C481-ghavavu}



\begin{quote}
Before the new one: say the Gujarati for to get cut, then the Gujarati for to heal.
\end{quote}

\subsection*{You'll want to know: \gu{ઘવાવું}}
\textbf{\gu{ઘવાવું}} — \emph{ghavāvũ} — \textquotedblleft{}to be wounded\textquotedblright{}.

\begin{quote}
\gu{ગાડી} \gu{અથડાઈ} \gu{ત્યારે} \gu{ડ્રાઈવર} \gu{ઘવાયો}.
\end{quote}

\begin{quote}
\emph{gāḍī athḍāī tyāre ḍrāīvar ghavāyo.}
\end{quote}

\begin{quote}
When the car crashed, the driver was injured.
\end{quote}

Past (he …): \gu{તે} \gu{ઘવાયો}, \emph{te ghavāyo}.

\begin{cousinweb}
\textbf{\gu{ઘવાવું}} (\emph{ghavāvũ}) is built on \textbf{\gu{ઘા}} (\emph{ghā}), \textquotedblleft{}a wound\textquotedblright{}, which goes back to Sanskrit \emph{ghāta}, \textquotedblleft{}blow, bruise\textquotedblright{} (CDIAL 4460), with the verb ending \emph{-vāvũ}.
\end{cousinweb}

\begin{grammarlens}[title={What you've built}]
One more word for describing what you lived and felt.
\end{grammarlens}

\subsection*{Your turn}
\begin{itemize}
\raggedright
  \item \emph{Say it:} \emph{ghavāvũ}
  \item \emph{Say it:} \emph{ghavāvũ}, once more
  \item \emph{Say it:} \emph{gāḍī athḍāī tyāre ḍrāīvar ghavāyo.} — When the car crashed, the driver was injured.
  \item \emph{Recall:} say the Gujarati for to get cut, then the Gujarati for to heal, then say \emph{ghavāvũ} again
\end{itemize}

\begin{tcolorbox}[breakable,colback=teal!4,colframe=teal!35!black,title={Before you move on}]
What does \gu{ઘવાવું} mean? (\textquotedblleft{}To be wounded\textquotedblright{}.) Say it once more.
\end{tcolorbox}
\section[\texorpdfstring{\gu{ઝરવું} (jharvũ)}{ઝરવું (jharvũ)}]{\gu{ઝરવું} (jharvũ) — to ooze; to trickle}

\label{lesson:GU-C481-jharvu}



\begin{quote}
Before the new one: say the Gujarati for to heal, then the Gujarati for to be wounded.
\end{quote}

\subsection*{You'll want to know: \gu{ઝરવું}}
\textbf{\gu{ઝરવું}} — \emph{jharvũ} — \textquotedblleft{}to ooze; to trickle\textquotedblright{}.

\begin{quote}
\gu{ઘામાંથી} \gu{હજી} \gu{લોહી} \gu{ઝરે} \gu{છે}.
\end{quote}

\begin{quote}
\emph{ghāmā̃thī hajī lohī jhare che.}
\end{quote}

\begin{quote}
Blood is still oozing from the wound.
\end{quote}

Past (it …): \gu{તે} \gu{લોહી} \gu{ધીમે} \gu{ઝર્યું}, \emph{te lohī dhīme jharyũ}.

\begin{cousinweb}
\textbf{\gu{ઝરવું}} (\emph{jharvũ}) goes back to a reconstructed \emph{*jharati}; CDIAL lists the Gujarati verb, \textquotedblleft{}to ooze, drip, shed, melt\textquotedblright{} (CDIAL 5346).
\end{cousinweb}

\begin{grammarlens}[title={What you've built}]
One more word for describing what you lived and felt.
\end{grammarlens}

\subsection*{Your turn}
\begin{itemize}
\raggedright
  \item \emph{Say it:} \emph{jharvũ}
  \item \emph{Say it:} \emph{jharvũ}, once more
  \item \emph{Say it:} \emph{ghāmā̃thī hajī lohī jhare che.} — Blood is still oozing from the wound.
  \item \emph{Recall:} say the Gujarati for to heal, then the Gujarati for to be wounded, then say \emph{jharvũ} again
\end{itemize}

\begin{tcolorbox}[breakable,colback=teal!4,colframe=teal!35!black,title={Before you move on}]
What does \gu{ઝરવું} mean? (\textquotedblleft{}To ooze; to trickle\textquotedblright{}.) Say it once more.
\end{tcolorbox}
\section[\texorpdfstring{\gu{છોલાવું} (cholāvũ)}{છોલાવું (cholāvũ)}]{\gu{છોલાવું} (cholāvũ) — to get grazed}

\label{lesson:GU-C481-cholavu}



\begin{quote}
Before the new one: say the Gujarati for to be wounded, then the Gujarati for to ooze.
\end{quote}

\subsection*{You'll want to know: \gu{છોલાવું}}
\textbf{\gu{છોલાવું}} — \emph{cholāvũ} — \textquotedblleft{}to get grazed\textquotedblright{}.

\begin{quote}
\gu{અરુણ} \gu{પડી} \gu{ગયો} \gu{અને} \gu{તેનો} \gu{ઘૂંટણ} \gu{છોલાઈ} \gu{ગયો}.
\end{quote}

\begin{quote}
\emph{Aruṇ paḍī gayo ane teno ghū̃ṭaṇ cholāī gayo.}
\end{quote}

\begin{quote}
Arun fell and grazed his knee.
\end{quote}

Past (it …): \gu{તે} \gu{ઘૂંટણ} \gu{છોલાઈ} \gu{ગયો}, \emph{te ghū̃ṭaṇ cholāī gayo}.

\begin{cousinweb}
\textbf{\gu{છોલાવું}} (\emph{cholāvũ}) is the intransitive partner of \textbf{\gu{છોલવું}} (\emph{cholvũ}), \textquotedblleft{}to peel\textquotedblright{}, which goes back to a reconstructed \emph{*chōll-} (CDIAL 5073): skin that peels off.
\end{cousinweb}

\begin{grammarlens}[title={What you've built}]
One more word for describing what you lived and felt.
\end{grammarlens}

\subsection*{Your turn}
\begin{itemize}
\raggedright
  \item \emph{Say it:} \emph{cholāvũ}
  \item \emph{Say it:} \emph{cholāvũ}, once more
  \item \emph{Say it:} \emph{Aruṇ paḍī gayo ane teno ghū̃ṭaṇ cholāī gayo.} — Arun fell and grazed his knee.
  \item \emph{Recall:} say the Gujarati for to be wounded, then the Gujarati for to ooze, then say \emph{cholāvũ} again
\end{itemize}

\begin{tcolorbox}[breakable,colback=teal!4,colframe=teal!35!black,title={Before you move on}]
What does \gu{છોલાવું} mean? (\textquotedblleft{}To get grazed\textquotedblright{}.) Say it once more.
\end{tcolorbox}
